Phase-field fracture models implemented as Abaqus user elements (UEL/UMAT)
enable regularized crack evolution, but production use on non-uniform meshes
raises three coupled questions: whether a published staggered baseline can be
reproduced with defensible reaction-force evidence; whether offline
MISESERI-driven pre-refinement can recover peak and pre-peak response at lower
cost than a uniform production mesh; and whether nonmatching state transfer
supports a controlled pre-peak checkpoint without overstating restart or
parallel claims.

This thesis documents a staged Abaqus workflow built around the Molnar and
Gravouil staggered single-notch benchmark~\cite{molnar2017}, phase-field
implementation context from related UEL literature~\cite{msekh2015,diddige2025},
and MISESERI pre-refinement practice motivated by Pandey and
Kumar~\cite{pandey2025}. Uniform mesh roles were fixed after supervisor
decisions that accept H2-PUB for fine RF--U validation, H1 for production
comparisons, and H0 for development. Offline refinement then produced a
corrected locally refined mesh with 14.7\% fewer physical elements than H1,
close peak/pre-peak RF--U agreement, and measured wall-time and memory savings
under a fair four-thread comparison. Controlled analytical transfer, Abaqus
ingestion, serial continuation, and a bounded pre-peak nonmatching release hold
were demonstrated; an offline active-set correction at a later checkpoint was
shown to be KKT-admissible under frozen history.

The same evidence sets hard limits. Unrestricted post-peak RF--U convergence
and crack-path mesh independence are not supported. Fixed active-set
continuation under further loading is disproven in the tested segment. The
D3D-A1 corrected package is not an accepted mechanically equilibrated restart.
Threaded shared-state characterization (P3-T4) produced no threaded evidence;
MPI and hybrid safety remain unqualified. ABAQUSER equivalence remains
externally blocked. Validated online adaptive remeshing is not claimed.

Practically, the work identifies where pre-refinement and pre-peak state
transfer are usable, quantifies accuracy/cost benefits within those scopes, and
provides a decision tree that blocks unsupported post-peak, restart, and
parallel generalizations.
